% ============================================================
%  SICORRE — Sistema de Control de Registros para la
%  Restitución de Derechos de NNA
%  Documento de Diseño — Capítulo 1: Organización del Diseño
% ============================================================

\chapter{Organización del Diseño}

El presente capítulo establece el marco general que rige la elaboración de este documento
de diseño para el \textbf{Sistema de Control de Registros para la Restitución de Derechos
	de Niños, Niñas y Adolescentes (SICORRE)}. Se definen los lineamientos, la metodología,
la notación y el alcance adoptados, con el fin de garantizar coherencia, trazabilidad y
calidad a lo largo de todas las fases del ciclo de vida del proyecto.

% ------------------------------------------------------------
\section{Lineamientos de Diseño}
% ------------------------------------------------------------

El diseño de SICORRE está regido por un conjunto de principios y directrices que aseguran
que el sistema responda tanto a las necesidades operativas de la Fundación de Restitución
de Derechos de Niños, Niñas y Adolescentes (en adelante, \textit{la Fundación}) como a
los estándares técnicos y éticos exigidos por la naturaleza sensible de la información
que gestiona.

\subsection*{Principios rectores}

\begin{itemize}
	\item \textbf{Centrado en el usuario.} El sistema debe ser intuitivo y accesible para
	el personal operativo de la Fundación, quienes en su mayoría provienen de perfiles
	de trabajo social y psicología, no necesariamente técnicos.
	
	\item \textbf{Confidencialidad y protección de datos.} Dado que SICORRE gestiona
	expedientes, diagnósticos y datos personales de menores en situación de
	vulnerabilidad, todo el diseño debe contemplar controles de acceso basados en
	roles (RBAC), cifrado en tránsito (TLS/HTTPS).
	\item \textbf{Separación de responsabilidades.} La arquitectura cliente-servidor
	(React + FastAPI + PostgreSQL) impone una separación clara entre presentación,
	lógica de negocio y persistencia, lo que facilita el mantenimiento y la
	evolución independiente de cada capa.
	
	\item \textbf{Trazabilidad y auditoría.} Toda operación de creación, modificación o
	consulta de un expediente debe quedar registrada con el identificador del usuario
	responsable, la fecha-hora y la acción realizada, garantizando un historial
	completo de cada caso.
	
	\item \textbf{Escalabilidad y mantenibilidad.} El diseño modular permitirá incorporar
	nuevas funcionalidades (p.\,ej., módulos de reportes estadísticos, alertas
	automáticas) sin impactar los módulos existentes.
	
	\item \textbf{Disponibilidad y tolerancia a fallos.} Se definirán estrategias de
	respaldo periódico de la base de datos PostgreSQL y se documentarán los
	procedimientos de recuperación ante desastres.
\end{itemize}

\subsection*{Estándares y normas adoptados}

\begin{itemize}
	\item \textbf{REST / OpenAPI 3.x} para la definición y documentación de los
	endpoints expuestos por el backend FastAPI.
	\item \textbf{PEP\,8} como guía de estilo para el código Python del backend.
	\item \textbf{Airbnb Style Guide} como convención de estilo para el código
	JavaScript/React del frontend.
	\item \textbf{Conventional Commits} para el registro de cambios en el repositorio
	de control de versiones (Git).

\end{itemize}

\subsection*{Restricciones de diseño}

\begin{itemize}
	\item El sistema deberá operar en la infraestructura existente de la Fundación o en
	un entorno de nube de bajo costo, por lo que el diseño evitará dependencias de
	plataformas propietarias de alto costo o terceros, enfocandoce en un entorno donde solo 
	tenga acceso personal autorizado de la organización.
	\item El backend será desarrollado exclusivamente con \textbf{FastAPI} (Python $\geq$ 3.11)
	y el frontend con \textbf{React} (Node.js $\geq$ 20 LTS), utilizando
	\textbf{PostgreSQL 15} como motor de base de datos relacional.
	\item La comunicación entre el frontend y el backend se realizará únicamente a través
	de la API REST documentada; no se permitirán accesos directos a la base de datos
	desde el cliente, no se permitira acceso a la API por terceros.
	\item El sistema deberá soportar al menos 30 usuarios concurrentes en la fase inicial
	de despliegue.
\end{itemize}

% ------------------------------------------------------------
\section{Metodología o Procedimiento o Técnicas}
% ------------------------------------------------------------

\subsection*{Metodología de diseño}

El diseño de SICORRE adopta un enfoque \textbf{orientado a componentes} alineado con el
patrón arquitectónico \textbf{MVC (Modelo-Vista-Controlador)}, adaptado a la separación
física de capas que impone la arquitectura elegida:

\begin{itemize}
	\item \textbf{Modelo:} representado por los esquemas de base de datos en PostgreSQL y
	los modelos ORM definidos con \textit{SQLAlchemy} en el backend.
	\item \textbf{Vista:} implementada como una Single Page Application (SPA) en
	\textit{React}, que consume los recursos expuestos por la API REST.
	\item \textbf{Controlador:} materializado en los \textit{routers} y \textit{servicios}
	de FastAPI, que encapsulan la lógica de negocio y orquestan el acceso a los datos.
\end{itemize}

El proceso de diseño sigue un ciclo iterativo-incremental, en el cual cada módulo
funcional (gestión de expedientes, diagnósticos, usuarios, reportes) es diseñado,
prototipado, revisado y refinado antes de pasar a la fase de construcción.

\subsection*{Técnicas empleadas}

\begin{itemize}
	\item \textbf{Modelado de datos relacional} mediante diagramas Entidad-Relación (ER)
	para la definición del esquema PostgreSQL.
	\item \textbf{Diseño de API-first:} los contratos de la API REST son definidos en
	OpenAPI antes de implementar los endpoints, facilitando el desarrollo paralelo
	del frontend y el backend.
	\item \textbf{Diagramas de casos de uso UML} para capturar las interacciones entre
	actores (trabajador social, coordinador, administrador del sistema) y el sistema.
	\item \textbf{Diagramas de secuencia UML} para modelar los flujos de operación
	críticos, como el registro de un nuevo expediente o la actualización de un
	diagnóstico.
	\item \textbf{Wireframes y prototipos de baja fidelidad} para validar la experiencia
	de usuario antes de la implementación de la interfaz en React o por su parte HTML antes de desarrollo y aprobación.
\end{itemize}

\subsection*{Herramientas}

\begin{table}[h!]
	\centering
	\caption{Herramientas utilizadas en el diseño de SICORRE}
	\begin{tabular}{|l|l|l|}
		\hline
		\textbf{Categoría}        & \textbf{Herramienta}          & \textbf{Propósito}                          \\ \hline
		Modelado UML              & draw.io / PlantUML            & Diagramas de casos de uso, secuencia, etc.  \\ \hline
		Diseño de API             & Swagger UI / OpenAPI 3.x      & Especificación y documentación de endpoints \\ \hline
		Modelado de BD            & dbdiagram.io / pgAdmin        & Diseño del esquema relacional               \\ \hline
		Prototipado UI            & Figma / HTML 				  & Wireframes y flujos de navegación           \\ \hline
		Control de versiones      & Git / GitHub                  & Gestión del código fuente                   \\ \hline
		Gestión del proyecto      & Teams / whatsapp              & Seguimiento de tareas e iteraciones         \\ \hline
		Documentación             & \LaTeX                        & Elaboración del presente documento          \\ \hline
	\end{tabular}
	\label{tab:herramientas}
\end{table}

% ------------------------------------------------------------
\section{Notación}
% ------------------------------------------------------------

La totalidad de los diagramas incluidos en este documento de diseño se elaboran conforme
al estándar \textbf{UML 2.x} (\textit{Unified Modeling Language}, versión 2 o superior). A continuación se describe
brevemente cada tipo de diagrama empleado y el propósito que cumple dentro del documento.

\begin{table}[h!]
	\centering
	\caption{Tipos de diagramas UML utilizados en SICORRE}
	\begin{tabular}{|p{4cm}|p{4cm}|p{5.5cm}|}
		\hline
		\textbf{Tipo de diagrama}      & \textbf{Notación principal}               & \textbf{Uso en este documento}                          \\ \hline
		Casos de uso                   & Actor, caso de uso (elipse), relaciones   & Describir interacciones entre usuarios y el sistema     \\ \hline
		Clases                         & Clases, atributos, métodos, asociaciones  & Modelar entidades del dominio y sus relaciones          \\ \hline
		Secuencia                      & Líneas de vida, mensajes, activaciones    & Detallar flujos de operaciones y llamadas entre capas   \\ \hline
		Componentes                    & Componentes, interfaces, dependencias     & Representar la arquitectura lógica del sistema          \\ \hline
		Despliegue                     & Nodos, artefactos, comunicaciones         & Describir la infraestructura física de ejecución        \\ \hline
		Entidad-Relación (ER)          & Entidades, atributos, relaciones (crow's foot) & Modelar el esquema de la base de datos PostgreSQL  \\ \hline
	\end{tabular}
	\label{tab:notacion}
\end{table}

Adicionalmente, para la documentación de los endpoints de la API REST se utiliza la
notación \textbf{OpenAPI 3.x}, que define de forma estandarizada los recursos, métodos
HTTP, parámetros, cuerpos de solicitud y códigos de respuesta.



% ------------------------------------------------------------
\section{Alcance}
% ------------------------------------------------------------

\subsection*{Descripción general}

El presente documento de diseño cubre el análisis y especificación arquitectónica del
\textbf{Sistema de Control de Registros para la Restitución de Derechos de Niños, Niñas
	y Adolescentes (SICORRE)}, desarrollado para la Fundación de Restitución de Derechos de
Niños, Niñas y Adolescentes en estado de orfandad debido al feminicidio en México.

SICORRE tiene como propósito digitalizar y estructurar los procesos actualmente gestionados
en papel, tales como la creación y seguimiento de expedientes, el registro de diagnósticos,
la gestión de fechas y citas, y el control de archivos relacionados con cada caso de NNA.

\subsection*{Módulos y funcionalidades dentro del alcance}

Los siguientes subsistemas y módulos son cubiertos por este documento de diseño:

\begin{enumerate}
	\item \textbf{Gestión de usuarios y roles:} administración de cuentas de acceso,
	asignación de roles (administrador, coordinador, trabajador social) y control
	de permisos.
	\item \textbf{Gestión de equipos multidisciplnarios} administración de equipos compuestos
	por diferentes usuarios y roles previamente registrados para dar seguimiento a los diagnosticos y
	planes para la resitución de derechos.
	\item \textbf{Gestión y registro de tutores} : administración y nuevos registros para tutores 
	con sus respectivos datos de contacto y asociación con los nnas.
	\item \textbf{Gestión y registro de nnas } : Gestión y registro principal de niños niñas y 
	adolescentes en el sistema.
	\item \textbf{Modulo de Restitución de derechos} : Modulo utilizado para consultar programas y actores en 
	materia de derechos que sirvan de ayuda para la restitución de los mismos para los nnas.
	\item \textbf{Consulta de catalogos} : consultas de catalogos establecidos previamente por la organización como 
	roles existentes, enfermedades, discapacidades, idiomas, domicilios (SEPOMEX). 
	\item \textbf{API REST:} especificación completa de los servicios que expone el
	backend FastAPI para consumo del frontend React.
	\item \textbf{Esquema de base de datos:} diseño del modelo relacional en PostgreSQL
	que soporta todos los módulos anteriores.
\end{enumerate}

\subsection*{Elementos fuera del alcance}

Los siguientes aspectos \textbf{no} son cubiertos por el presente documento de diseño:

\begin{itemize}
	\item Integración con sistemas externos de gobierno (SIPINNA, DIF, Fiscalía) en la
	fase inicial; podrá contemplarse en versiones futuras.
	\item Módulo de comunicación directa con familias o tutores (mensajería, portal
	ciudadano).
	\item Aplicación móvil nativa; la solución es exclusivamente web responsiva.
	\item Módulo de firma electrónica avanzada o firma digital con valor legal.
	\item Infraestructura de despliegue en producción (configuración de servidores,
	balanceo de carga, DNS).
	\item Migración de expedientes físicos históricos en papel; este proceso se realizará
	de forma manual y operativa fuera del alcance del sistema o por asistencia de modelos de IA propuestos 
	hechos.
\end{itemize}

\bigskip
\noindent

La Tabla~\ref{tab:alcance-resumen} presenta un resumen del alcance por capa de la
arquitectura.

\begin{table}[h!]
	\centering
	\caption{Resumen del alcance por capa arquitectónica}
	\begin{tabular}{|l|l|c|}
		\hline
		\textbf{Capa}         & \textbf{Tecnología}         & \textbf{Cubierta en este doc.} \\ \hline
		Presentación (UI)     & React (Node.js)             & \checkmark                     \\ \hline
		Lógica de negocio     & FastAPI (Python)            & \checkmark                     \\ \hline
		Persistencia          & PostgreSQL 15               & \checkmark                     \\ \hline
		Infraestructura       & Servidor / Nube             & \texttimes                     \\ \hline
		Integraciones externas & APIs gubernamentales       & \texttimes                     \\ \hline
	\end{tabular}
	\label{tab:alcance-resumen}
\end{table}